\clearpage
\SourceLeaf{28}
\begin{center}\large [FRAGMENTS ON THE THEORY OF ALGEBRAIC INTEGERS FORMED FROM A CUBE ROOT.]\end{center}

\begin{center}[1.]\end{center}\par\noindent\hspace*{\parindent}Let
\[
a+b\nu+c\nu\nu = t, \qquad \nu^3 = n,
\]
\noindent and set
\[
a^3+nb^3+nnc^3-3nabc = \varphi t,
\]
\[
(bb-ac)\nu\nu+(ncc-ab)\nu+(aa-nbc) = Ft.
\]
\noindent\textls{Theorem I.}
\[
tFt = \varphi t,
\]
\[
\left.\begin{aligned}
\varphi kt &= k^3 \varphi t \\
Fkt &= kk\,Ft
\end{aligned}\right\}\ \text{if } k \text{ is a number}
\]
\[
\varphi tu = \varphi t . \varphi u, \qquad Ftu = Ft . Fu,
\]
\[
\varphi(Ft) = (\varphi t)^2, \qquad F(Ft) = (\varphi t)\,t.
\]

If the coefficients in $t$ have no common divisor, while those in $Ft$ have a
common divisor $k$, and we set $Ft = kt'$, then
\[
F(Ft) = kk\,Ft' = t\varphi t
\]%
